\section{Error correcting codes}\label{sec:error_correcting_codes}

\paragraph{Noisy channels}

\begin{concept}\label{con:communication_channel}\mcite[19]{Shannon1948TheoryOfCommunication}
  Information is transmitted over a \term[ru=канал связи (\cite[258]{Яблонский2003ДискретнаяМатематика}), en=channel of communication (\cite[879]{GowersEtAl2008PrincetonCompanion})]{communication channel}, which can range from spoken word to public computer networks.
\end{concept}
\begin{comments}
  \item Mathematically, a channel can be described by different formalisms, depending on the problem at hand.
  \begin{itemize}
    \item Noisy channels, of interest for error-correcting codes, are defined in \cref{def:noisy_channel} as functions with an auxiliary noise parameter.

    \item The distinction between \hyperref[con:privacy]{secure and insecure} channels is important in \fullref{sec:cryptography}, but the channels themselves are of little theoretical interest.

    \item How the throughput of multiple channels is combined can be studied via \hyperref[def:directed_graph]{directed graphs}. In this case only the capacity of a channel is relevant.
  \end{itemize}
\end{comments}

\begin{definition}\label{def:noisy_channel}\mcite[19]{Shannon1948TheoryOfCommunication}
  A \hyperref[con:communication_channel]{communication channel} can introduce \term[ru=помехи (\cite[288]{Яблонский2003ДискретнаяМатематика})]{noise}. Formally, this can be described by some function \( f: \mscrS \times \mscrN \to \mscrE \), where \( \mscrS \) and \( \mscrE \) are sent and received messages, while noise is modeled by the set \( \mscrN \) consists of \hyperref[def:random_variable]{random variables}.

  \begin{thmenum}
    \thmitem{def:communication_channel/symmetric}\mimprovised Suppose that \( \mscrS = \mscrE = \set{ 0, 1 }^* \), i.e. the sent and received messages are \hyperref[def:bit_string]{bit strings}.

    Fix a probability \( p \in (0, 1) \) and let \( \mscrN \) be a set of sequences \( N = \seq{ N_i }_{i=1}^\infty \) of independent Bernoulli random variables with probability \( p \).

    We define the \term[en=binary symmetric channel (\cite[75]{ConwaySloane1998SpherePackings})]{binary symmetric channel} with parameter \( p \) as
    \begin{equation*}
      f(w_1 \ldots w_m, N) = (w_1 \oplus N_1) \ldots (w_m \oplus N_m).
    \end{equation*}

    This channel can alter any bit of \( w = w_1 \ldots w_m \) with probability \( p \).

    We can consider a more general \hyperref[def:formal_language/alphabet]{alphabet} \( \Sigma \) than \( \set{ 0, 1 } \). We must then alter the noise so that \( N = \seq{ (N_i, A_i) }_{i=1}^\infty \), where \( N_i \) is as before, while \( A_i \) is a uniformly random element of \( \Sigma \). We then replace \( w_i \) with \( A_i \) if \( N_i = 1 \) and preserve it otherwise.

    We call this the \( \Sigma \)-\term[en=symmetric channel (\cite[879]{GowersEtAl2008PrincetonCompanion})]{symmetric channel}.
  \end{thmenum}
\end{definition}

\begin{definition}\label{def:error_correcting_code}\mcite[75]{ConwaySloane1998SpherePackings}
  An \term[ru=самокорректирующейся код (\cite[289]{Яблонский2003ДискретнаяМатематика}), en=error-correcting code (\cite[def. 19.9]{AroraBarak2009ComputationalComplexity})]{error-correcting code} of \term{length} \( n > 0 \) over the \hyperref[def:finite_field]{finite field} \( \BbbF_q \) is simply a subset of \( \BbbF_q^n \)

  \begin{thmenum}
    \thmitem{def:error_correcting_code/codeword} We call the elements of a code \term[ru=коды (\cite[258]{Яблонский2003ДискретнаяМатематика})]{codewords}.

    \thmitem{def:error_correcting_code/distance} We define the \term{minimal distance} of a code as the smallest \hyperref[def:hamming_distance]{Hamming distance} between (distinct) codewords.

    \thmitem{def:error_correcting_code/linear}\mcite[75]{ConwaySloane1998SpherePackings} If the code is a \hyperref[def:vector_space/subobject]{vector subspace} of \( \BbbF_q^n \), we say that it is \term{linear}.
  \end{thmenum}
\end{definition}
\begin{comments}
  \item We will use prefixes like \enquote{binary} or \enquote{ternary} for the code itself, based on how we call the corresponding \hyperref[def:positional_number_system]{positional number systems}.

  \item Some authors like \textcite[75]{ConwaySloane1998SpherePackings} and \textcite[2]{Berlekamp2015AlgebraicCodingTheory} shorten \enquote{error-correcting code} to \enquote{code}. We avoid doing so due to ambiguity.

  \item \Textcite[def. 19.9]{AroraBarak2009ComputationalComplexity} instead define an error-correcting code as a function \( f: \Sigma^m \to \Sigma^n \) for some alphabet \( \Sigma \). This approach has its appeal since codes are implicitly connected with their encoding algorithms, however we follow the more established practice of treating a code as a set.
\end{comments}

\paragraph{Parity check}

\begin{algorithm}[Parity check code]\label{alg:parity_check_code}\mcite[example 9.3]{LidlNiederreiter1997FiniteFields}
  A simple binary \hyperref[def:error_correcting_code]{error-correcting code} is the \term{parity check code}, which extends the message \( w = w_n \cdots w_m \) with a single bit
  \begin{equation*}
    w_{m+1} \coloneqq \bigoplus_{k=1}^m w_k.
  \end{equation*}
\end{algorithm}
\begin{comments}
  \item This algorithm is implemented in \cite{notebook:code} as
  \begin{CodeRefDisplay}
    \GetCodeRef{alg:parity_check_code}
  \end{CodeRefDisplay}
\end{comments}

\begin{proposition}\label{thm:parity_check_codeword_characterization}
  A \hyperref[def:bit_string]{bit string} of length at least \( 2 \) is a codeword for \fullref{alg:parity_check_code} if and only if it has an even number of ones.
\end{proposition}
\begin{comments}
  \item If the string is a \hyperref[def:ring_of_unsigned_integers]{binary encoding} of some number, the condition is equivalent for the number to be even.
\end{comments}
\begin{proof}
  Denote the number of occurrence of \( 1 \) in a bit string \( s \) by \( N(s) \). Clearly \( N(s) \) is even if and only if the sum of its digits is \( 0 \).

  Fix some string \( c = w w_{m+1} \), where \( w = w_1 \ldots w_m \). It is a codeword if and only if
  \begin{equation*}
    w_{m+1} = \bigoplus_{k=1}^m w_k.
  \end{equation*}

  We have two cases:
  \begin{itemize}
    \item In case \( N(w) \) is even, \( w_{m+1} = 0 \) if and only if \( N(w w_{m+1}) = N(w) \) is even.
    \item In case \( N(w) \) is odd, \( w_{m+1} = 1 \) if and only if \( N(w w_{m+1}) = N(w) + 1 \) is even.
  \end{itemize}

  Therefore, \( w_{m+1} \) is a valid parity check digit if and only if \( N(w w_{m+1}) \) is even.
\end{proof}

\begin{corollary}\label{thm:parity_check_distance}
  \Fullref{alg:parity_check_code} has \hyperref[def:error_correcting_code/distance]{minimal distance} \( 2 \).
\end{corollary}
\begin{proof}
  \Cref{thm:parity_check_codeword_characterization} characterizes the codewords as those with an even number of ones. Therefore, two codewords must differ in at least two positions.
\end{proof}

\begin{proposition}\label{thm:parity_check_linearity}
  \Fullref{alg:parity_check_code} is \hyperref[def:error_correcting_code/linear]{linear}.
\end{proposition}
\begin{proof}
  \begin{equation*}
    \begin{pmatrix}
      w_1 \\
      w_2 \\
      \vdots \\
      w_m \\
      w_{m+1}
    \end{pmatrix}
    =
    \begin{pmatrix}
      1      & 0      & \cdots & 0 \\
      0      & 1      & \ddots & 0 \\
      \vdots & \vdots & \ddots & 0 \\
      0      & 0      & \cdots & 1 \\
      1      & 1      & \cdots & 1
    \end{pmatrix}
    \cdot
    \begin{pmatrix}
      w_1 \\
      w_2 \\
      \vdots \\
      w_m
    \end{pmatrix}
  \end{equation*}
\end{proof}

\paragraph{Reed-Solomon codes}

\begin{algorithm}[Reed-Solomon code]\label{alg:reed_solomon_code}
  We will now describe the \hyperref[def:error_correcting_code]{error-correcting} \term{Reed-Solomon code} over the \hyperref[def:finite_field]{finite field} \( \BbbF_q \).

  Fix positive integers \( m \leq n \leq q \) and \( n \) distinct \emph{nonzero} elements \( \alpha_1, \ldots, \alpha_n \) of \( \BbbF_q \).

  For a given message \( w = w_1 \cdots w_m \), we first define the polynomial
  \begin{equation*}
    f(X) \coloneqq \sum_{k=0}^{m-1} w_k X^k
  \end{equation*}
  and then the codeword as the string
  \begin{equation*}
    f(\alpha_1) \cdots f(\alpha_n).
  \end{equation*}
\end{algorithm}
\begin{comments}
  \item This algorithm is implemented in \cite{notebook:code} as
  \begin{CodeRefDisplay}
    \GetCodeRef{alg:reed_solomon_code}
  \end{CodeRefDisplay}
\end{comments}

\begin{proposition}\label{thm:reed_solomon_distance}
  The \hyperref[alg:reed_solomon_code]{Reed-Solomon code} has \hyperref[def:error_correcting_code/distance]{minimal distance} \( n - m + 1 \).
\end{proposition}
\begin{proof}
  For the codewords, \( v \neq w \) consider the difference of their polynomials
  \begin{equation*}
    \sum_{k=0}^{m-1} (v_k - w_k) X^k.
  \end{equation*}

  It has degree at most \( m - 1 \), so it has at most \( m - 1 \) roots according to \cref{thm:polynomial_degree_bounds_root_count}. Therefore, the codewords for \( v \) and \( w \) differ in at least \( n - m + 1 \) positions.
\end{proof}

\begin{proposition}\label{thm:reed_solomon_linearity}
  The \hyperref[alg:reed_solomon_code]{Reed-Solomon code} is \hyperref[def:error_correcting_code/linear]{linear}.
\end{proposition}
\begin{proof}
  \begin{equation*}
    \begin{pmatrix}
      f(\alpha_1) \\
      f(\alpha_2) \\
      \vdots \\
      f(\alpha_n)
    \end{pmatrix}
    =
    \begin{pmatrix}
      1      & \alpha_1 & \cdots & \alpha_1^{m-1} \\
      \vdots & \vdots   & \ddots & \vdots \\
      1      & \alpha_n & \cdots & \alpha_n^{m-1}
    \end{pmatrix}
    \cdot
    \begin{pmatrix}
      w_1 \\
      w_2 \\
      \vdots \\
      w_m
    \end{pmatrix}
  \end{equation*}
\end{proof}
